\section*{Chapitre \ref{ch:g1:materials-around-us} --- De quoi sont faits les objets}

\begin{solution}{exo:g1:materials-around-us:1}
Une chaise est un \omterm{def:g1:materials-around-us:object}{objet}. Le bois est un \omterm{def:g1:materials-around-us:material}{matériau}~: on peut en faire une
chaise.
\end{solution}

\begin{solution}{exo:g1:materials-around-us:2}
Vitre~: verre. Clou~: métal. Galet~: pierre. Livre~: papier. Chaussette~:
tissu (laine ou coton).
\end{solution}

\begin{solution}{exo:g1:materials-around-us:3}
Bois~: un crayon, une porte (ou une chaise, une table). Plastique~: un
bol, un seau (ou un jouet, une bouteille).
\end{solution}

\begin{solution}{exo:g1:materials-around-us:4}
La fourchette en métal~: c'est la seule qui ne soit ni en bois ni en
plastique. (Deux \omterm{def:g1:materials-around-us:object}{objets} sont en bois et un en plastique.)
\end{solution}

\begin{solution}{exo:g1:materials-around-us:5}
Le verre~: il est transparent.
\end{solution}

\begin{solution}{exo:g1:materials-around-us:6}
Le clou en acier. Un aimant attire le fer et l'acier, pas le verre ni le
bois.
\end{solution}

\begin{solution}{exo:g1:materials-around-us:7}
Deux \omterm{def:g1:materials-around-us:object}{objets} flottent (le bloc de bois et le bouchon en plastique)~; trois
coulent (le galet, la bille et le clou).
\end{solution}

\begin{solution}{exo:g1:materials-around-us:8}
Les lames sont en métal, les poignées en plastique. Les ciseaux sont
faits de deux matériaux.
\end{solution}

\begin{solution}{exo:g1:materials-around-us:9}
Dans le cerceau bleu~: elle est en métal. Elle n'est pas dans le cerceau
orange parce qu'un aimant ne l'attire pas~: elle est en aluminium, pas
en fer ni en acier.
\end{solution}

\begin{solution}{exo:g1:materials-around-us:10}
Une fenêtre doit laisser entrer la lumière, et on doit voir à travers.
Le verre est transparent~; le bois ne l'est pas.
\end{solution}

\begin{solution}{exo:g1:materials-around-us:11}
«~Oui~»~: le clou, le trombone et la clé, 3 \omterm{def:g1:materials-around-us:object}{objets}. «~Non~»~: le galet,
la cuillère, la canette, la bille et la chaussette, 5 \omterm{def:g1:materials-around-us:object}{objets}.
$3 + 5 = 8$~: chaque \omterm{def:g1:materials-around-us:object}{objet} a sa place.
\end{solution}
